% Tapis de Quillen, with Quillen's offprint and preprint (cote 162-5, typed part): archive cote 162-5, site page 43 (PDF page 44).
% typescript. Draft compared: none.
% AI-transcribed by Claude Opus 5.5 (claude-opus-5-5, adaptive thinking, effort high), Message Batches API, batch t162_full (msgbatch_01UAh3zyWPd454AeDkhnmQ33), 2026-11-13; not yet checked by eye.
\scanpage{162-5/43}
III   Point de vue "motivique" en théorie du cobordisme.

Soit C la catégorie des variétés différentiables (pas nécessairement
orientables), \inserted[typed]{les morphismes étant les classes d'homotopie d'applications continues.} Si B est une catégorie, on s'intéresse aux couples
$(F_{.},F^{\cdot})$ d'un foncteur covariant et d'un foncteur contravariant de C
dans B, satisfaisant les conditions que pour tout X \hand{$\in$} Ob C, on a $F_{.}(X)=$
$F^{\cdot}(X)$, et que si on a un produit fibré ordinaire
\diagram{
\begin{tikzcd}[column sep=small, row sep=small]
 & X' \arrow[dl, "g'"'] \arrow[dr, "f'"] & \\
X \arrow[dr, "f"'] & & Y' \arrow[dl, "g"] \\
 & Y &
\end{tikzcd}
}
,
avec f et g transversaux, alors on a $g^{\cdot}f_{.} = f!\ g'^{\cdot}$ \inserted[typed]{(on pose $f_{.}=F_{.}(f),f^{\cdot}=F^{\cdot}(f)$)}. La question est
de trouver une B et un couple $(F_{.},F^{\cdot})$ qui soient "universels" dans un
sens évident. Quillen prouve que la construction suivante donne une
telle solution.

Pour toute variété X\hand{$\in$} Ob C , on définit un anneau commutatif
B(X) (l'anneau de bordisme de X) en prenant comme éléments les classes
\xout{de} d'éléments de C au dessus de X, pour la relation de cobordisme:
$f':X'\to X$ et $f'':X''\to X$ sont équivalents, s'il existe une variété à bord
Z, et un $f:Z\to X$, tel que bord $Z = X'$\hand{$\amalg$}$X''$ et f' et f" sont induits
par f. La somme de B(X) est la somme disjointe (c'est une loi de groupe,
\xout{l'inverse étant} où tout élément est d'ordre \deleted{\unsure{1}}\hand{2} \deleted{\unsure{ou 2}}), le produit est
donné par le produit fibré, en utilisant le théorème de trahsversalité
de Thom pour bouger un peu f", de sorte à le rendre transversal à f'.
\inserted{\uline{NB} L'anneau est de plus \unsure{gradué} par la dimension relative ...}
Si on a $f: X\to Y$, on définit $f_{.} : B(X)\to B(Y)$ de façon évidente (toute
variété sur X est sur Y), et $f^{\cdot}: B(Y)\to B(X)$ par produit fibré (utili-%
sant encore le th. de transversalité). De cette façon, B(X) devient un
foncteur-anneau \inserted{(gradué)} contravariant en X, et un foncteur groupe covariant en

\mspage{43}
